\documentclass[12pt]{article}
\usepackage{amsfonts,amsthm,amsmath,amssymb,mathtools}
\usepackage{mathrsfs}
\usepackage{enumitem}
\usepackage{tikz-cd}
\usepackage{geometry}
\usepackage{mlmodern}
\usepackage{graphicx}

\geometry{letterpaper, portrait, margin=1in}

\setcounter{section}{0}

\renewcommand{\thesection}{\S\arabic{section}}
\renewcommand{\thesubsection}{\arabic{section}}
\DeclareMathOperator{\im}{im}
\DeclareMathOperator{\id}{id}

\newtheorem{thm}{Theorem}
\newenvironment{theorem}[1]{
	\renewcommand{\thethm}{#1}
	\begin{thm}
		}{\end{thm}}

\newtheorem{lma}{Lemma}
\newenvironment{lemma}[1]{
	\renewcommand{\thelma}{#1}
	\begin{lma}
		}{\end{lma}}

\theoremstyle{definition}
\newtheorem{ex}{Assignment}
\newenvironment{exercise}[1]{
	\renewcommand{\theex}{#1}
	\begin{ex}
		}{\end{ex}}

\title{MTH 732 Topology - Homework 10}
\author{Connor Mooneyhan}
\date{April 23, 2026}

\begin{document}

\maketitle

\begin{ex}
	Show that the degree of the antipodal map $A : S^n \to S^n$ given by $A(x) = -x$ is $\deg(A) = (-1)^{n+1}$.
\end{ex}
\begin{proof}
  First notice that if we only ``flip'' one of the coordinates with a map $f$, then $f$ has degree -1.
  Composing maps $S^n \to S^n$ causes the degrees to multiply, so if we consider $A$ to be the composition of $n+1$ single-coordinate flips (letting $S^n$ sit inside $R^{n+1}$), then we get that $\deg(A) = (-1)^{n+1}$ as desired.
\end{proof}

\begin{ex}
	Describe a CW complex structure on the 3-torus $T^3 = S^1\times S^1 \times S^1$ and compute its homology.
\end{ex}
\begin{proof}
	Using the product CW structure, we can see that since $S^1 \times S^1$ has one 0-cell, two 1-cells, and one 2-cell and $S^1$ has one 0-cell and one 1-cell, the product will give one 0-cell, three 1-cells, three 2-cells, and one 3-cell.
	The 1-cells are attached by the only possible map, namely the constant map, so $\partial_1$ is the zero map.
	The 2-cells are attached similar to on the torus, where each 1-cell is traversed twice, once backward and once forward, so $\partial_2$ is also the zero map.
	Finally, the 3-cell connects all of the 2-cells, and once again in a way that traverses each one once in each orientation, so $\partial_3$ is also zero.
	Thus we can read off the homology straight from the chain groups and we get \begin{equation*}
		H_n(T^3) = \begin{cases}
			\mathbb{Z}   & n = 0, 3         \\
			\mathbb{Z}^3 & n = 1, 2         \\
			0            & \text{otherwise}.
		\end{cases}
	\end{equation*}
\end{proof}

\begin{ex}
	For $n < m$, find the homology groups of the quotient space $\mathbb{RP}^m/\mathbb{RP}^n$, where $\mathbb{RP}^n$ sits inside $\mathbb{RP}^m$ as its $n$-skeleton.
\end{ex}
\begin{proof}
	First, note that $\mathbb{RP}^m$ has one cell in each dimension from 0 through $m$, so $\mathbb{RP}^m/\mathbb{RP}^n$ has one cell in each dimension from $n + 1$ through $m$, \textit{and} one 0-cell (which arises from the modded-out $\mathbb{RP}^n$).
	So the cellular chain complex looks like this: \begin{equation*}
		\begin{tikzcd}
			\cdots \arrow[r, "\partial_{m+2}"] & 0 \arrow[r, "\partial_{m+1}"] & \mathbb{Z} \arrow[r, "\partial_m"] & \cdots \arrow[r, "\partial_{n+2}"] & \mathbb{Z} \arrow[r, "\partial_{n+1}"] & 0 \arrow[r, "\partial_n"] & \cdots \arrow[r, "\partial_2"] & 0 \arrow[r, "\partial_1"] & \mathbb{Z} \arrow[r, "\partial_0"] & 0
		\end{tikzcd},
	\end{equation*}
	where the boundary maps are all 0 since the attaching maps are all constant maps to the 0-cell.
	Given this, we can read off \begin{equation*}
		H_k\left(\mathbb{RP}^m/\mathbb{RP}^n\right) = \begin{cases}
			\mathbb{Z} & k = 1               \\
			\mathbb{Z} & n + 1 \leq k \leq m \\
			0          & \text{otherwise}.
		\end{cases}
	\end{equation*}
\end{proof}

\begin{ex}
	In class, we saw that the real Grassmannian $\mathrm{Gr}(2, 4)$ has a CW complex structure coming from all possible reduced row echelon forms of $2 \times 4$ matrices.
	In particular, there are exactly two 2-cells: \begin{equation*}
		e^2_1 = \begin{pmatrix}0 & 1 & 0 & x \\ 0 & 0 & 1 & y\end{pmatrix} \hspace{2em} \text{and} \hspace{2em} e^2_2 = \begin{pmatrix}1 & x & y & 0 \\ 0 & 0 & 0 & 1 \end{pmatrix}
	\end{equation*}
	where $(x, y) \in \mathbb{R}^2$.
	Show that the attaching maps $\varphi^2_i : \partial D^2_i \to X^1$, $i = 1, 2$, of these cells onto the 1-skeleton of $\mathrm{Gr}(2, 4)$ (which is a circle) have degree 2.
\end{ex}
\begin{proof}
  We are given the two 2-cells $e^2_1$ and $e^2_2$, and we also have the 1-cell and 0-cell \begin{equation*}
    e^1 = \begin{pmatrix}0 & 1 & x & 0 \\ 0 & 0 & 0 & 1\end{pmatrix} \hspace{2em} \text{and} \hspace{2em} e^0 = \begin{pmatrix} 0 & 0 & 1 & 0 \\ 0 & 0 & 0 & 1 \end{pmatrix}.
  \end{equation*}
  Since $e^0$, $e^1$, and $e^2_1$ each have zeroes in the first column, they form a copy of $\mathrm{Gr}(2, 3)$, which is isomorphic to $\mathrm{Gr}(1, 3)$, which is $\mathbb{RP}^2$.
  We already know that the boundary map of the 2-cell in $\mathbb{RP}^2$ has degree 2, so we're done $\varphi^2_1$.

  Now we notice that the second row of $e^2_2$ is constant, so that any plane represented by $e^2_2$ contains the vector $(0, 0, 0, 1)$.
  Thus taking $e^2_2$ together with $e^1$ and $e^0$, we have again the set of planes in a 3-dimensional subspace, and by the same argument above $\varphi^2_2$ has degree 2.
\end{proof}

\begin{ex}
	Show that any finite, path-connected CW complex $X$ is homotopy equivalent to a CW complex $Y$ having exactly \textbf{one} 0-cell.
\end{ex}
\begin{proof}
  Since $X$ is path-connected, there cannot be any isolated 0-cells, since that would produce multiple path components.
  Thus all the 0-cells have to be connected by some path of 1-cells in the 1-skeleton.
  We can choose, then, some spanning tree of the 1-skeleton and contract it to get a homotopy equivalent CW complex which has exactly one 0-cell.
  Note that the homotopy equivalence comes from the fact that the tree itself is contractible since it has no loops, so contracting it gives a homotopy equivalent complex.
\end{proof}

\begin{ex}
	If $p : \widetilde{X} \to X$ is a $k$-fold covering map (every preimage $p^{-1}(x)$ has $k$ elements) and $X$ is a finite CW complex, show that $\widetilde{X}$ has a natural CW complex structure coming from the lifts of the cells of $X$ (more precisely, lifting the characteristic maps $\phi^n_\alpha$ restricted to $\mathrm{int}(D^n_\alpha)$).
	Conclude that $\chi(\widetilde{X}) = k\chi(X)$.
\end{ex}
\begin{proof}
  For any $x$ in the interior of some $n$-cell $D^n_\alpha$, $p$ restricts to a homeomorphism on any component of $p^{-1}(\mathrm{int}(D^n_\alpha))$.
  Thus there are $k$ $n$-cells in $\widetilde{X}$ for each $n$-cell in $X$.
  Meanwhile, if $x$ is on the boundary of some $n$-cell, then it is also on the boundary of some $k$-cell for $k < n$ (for at least one such $k$, but possibly more), so that any neighborhood $U$ of $x$ intersects the interior of all such cells it is on the boundary of.
  Then the attaching maps pull back to attaching maps in $\widetilde{X}$, and we get that, for a choice of preimage component, we have an entire copy of $X$ of which it is a part.
  Thus, as desired, we get that $\chi(\widetilde{X}) = k\chi(X)$ since the number of cells in each dimension is multiplied by $k$, and then we factor $k$ out of the alternating sum.
\end{proof}

\begin{ex}
	Show that tensoring the short exact sequence \begin{tikzcd}[column sep=small]0 \arrow[r] & \mathbb{Z} \arrow[r, "n"] & \mathbb{Z} \arrow[r] & \mathbb{Z}_n \arrow[r] & 0\end{tikzcd} with an Abelian group $G$ can destroy exactness on the left (find such a group $G$).
\end{ex}
\begin{proof}
  Choose $G = \mathbb{Z}_n$.
  Then the sequence becomes \begin{tikzcd}[column sep=small]0 \arrow[r] & \mathbb{Z}_n \arrow[r, "n"] & \mathbb{Z}_n \arrow[r] & \mathbb{Z}_n \otimes \mathbb{Z}_n \arrow[r] & 0\end{tikzcd}.
    Since the multiplication by $n$ map sends 1 to 0, it is not injective, so the sequence is not left exact.
\end{proof}

\begin{ex}
	Show that $\mathbb{RP}^3$ and $\mathbb{RP}^2 \vee S^3$ are not homotopy equivalent.
\end{ex}
\begin{proof}
  I know this has something to do with the cup product, but I'm not getting there with this one.
\end{proof}

\end{document}
